% GENERATED FILE. Edit canonical lessons, then run npm run generate:books.
% canonical-source-hash: fnv1a64:503c36647e83d4c9
% canonical-lessons: FR-C291-gendarme, FR-C291-gendarmerie, FR-C291-patrouiller, FR-C291-patrouille, FR-C291-brigade, FR-R291-revision-284-291

\chapter{Les gendarmes}
\label{ch:fr-les-gendarmes}
\hlchaptermodality{291}

\begin{chapteropening}
\textbf{By the end of this chapter:} \emph{I can talk about the police on patrol: the gendarme, the gendarmerie, patrolling, a patrol and a squad.}

Complete the last lesson of chapter 291: I can talk about the police on patrol: the gendarme, the gendarmerie, patrolling, a patrol and a squad.
\end{chapteropening}

\section[le gendarme]{le gendarme — the gendarme}

\label{lesson:FR-C291-gendarme}



\begin{quote}
Before the new one: say the French for the offender, then the French for crime.
\end{quote}

\subsection*{You'll want to know: le gendarme}
\textbf{le gendarme} — \textquotedblleft{}the gendarme\textquotedblright{}.

\begin{quote}
Les gendarmes ont arrêté la voiture sur la route.
\end{quote}

\begin{quote}
\emph{The gendarmes stopped the car on the road.}
\end{quote}

In France, \emph{les gendarmes} are soldiers who do police work, mostly outside big towns; \emph{les policiers} work in the cities.

\begin{cousinweb}
\emph{gendarme} comes from \emph{gens d'armes} \textquotedblleft{}men of arms\textquotedblright{} (EtymDB lists the word without a parent): \emph{gens} goes back to Latin \emph{gēns} \textquotedblleft{}a people, a clan\textquotedblright{} and \emph{armes} to Latin \emph{arma} \textquotedblleft{}weapons\textquotedblright{} (Lewis \& Short).
\end{cousinweb}

\begin{grammarlens}[title={What you've built}]
One more word for telling, explaining and describing at B1.
\end{grammarlens}

\subsection*{Your turn}
\begin{itemize}
\raggedright
  \item \emph{Say it:} \emph{le gendarme}
  \item \emph{Say it:} \emph{le gendarme}, once more
  \item \emph{Say it:} a sentence of your own with \emph{le gendarme}
  \item \emph{Recall:} say the French for the offender, then the French for crime, then say \emph{le gendarme} again
\end{itemize}

\begin{tcolorbox}[breakable,colback=teal!4,colframe=teal!35!black,title={Before you move on}]
What does le gendarme mean? (\textquotedblleft{}The gendarme\textquotedblright{}.) Say it once more.
\end{tcolorbox}
\section[la gendarmerie]{la gendarmerie — the gendarmerie; its offices}

\label{lesson:FR-C291-gendarmerie}



\begin{quote}
Before the new one: say the French for crime, then the French for the gendarme.
\end{quote}

\subsection*{You'll want to know: la gendarmerie}
\textbf{la gendarmerie} — \textquotedblleft{}the gendarmerie; its offices\textquotedblright{}.

\begin{quote}
Elle a perdu son sac et elle a appelé la gendarmerie.
\end{quote}

\begin{quote}
\emph{She lost her bag and called the gendarmerie.}
\end{quote}

Both the force and the building.

\begin{cousinweb}
\emph{gendarmerie} is built on \emph{gendarme} (EtymDB), from \emph{gens d'armes} \textquotedblleft{}men of arms\textquotedblright{}.
\end{cousinweb}

\begin{grammarlens}[title={What you've built}]
One more word for telling, explaining and describing at B1.
\end{grammarlens}

\subsection*{Your turn}
\begin{itemize}
\raggedright
  \item \emph{Say it:} \emph{la gendarmerie}
  \item \emph{Say it:} \emph{la gendarmerie}, once more
  \item \emph{Say it:} a sentence of your own with \emph{la gendarmerie}
  \item \emph{Recall:} say the French for crime, then the French for the gendarme, then say \emph{la gendarmerie} again
\end{itemize}

\begin{tcolorbox}[breakable,colback=teal!4,colframe=teal!35!black,title={Before you move on}]
What does la gendarmerie mean? (\textquotedblleft{}The gendarmerie; its offices\textquotedblright{}.) Say it once more.
\end{tcolorbox}
\section[patrouiller]{patrouiller — to patrol}

\label{lesson:FR-C291-patrouiller}



\begin{quote}
Before the new one: say the French for the gendarme, then the French for the gendarmerie.
\end{quote}

\subsection*{You'll want to know: patrouiller}
\textbf{patrouiller} — \textquotedblleft{}to patrol\textquotedblright{}.

\begin{quote}
Des policiers patrouillent dans le parc la nuit.
\end{quote}

\begin{quote}
\emph{Police officers patrol the park at night.}
\end{quote}

\emph{Patrouiller dans un quartier} — to patrol a neighbourhood.

\begin{cousinweb}
\emph{patrouiller} comes from Old French \emph{patrouiller} \textquotedblleft{}to paddle about, to tramp in mud\textquotedblright{} (EtymDB), built on \emph{patte} \textquotedblleft{}paw\textquotedblright{}, whose origin is not settled.
\end{cousinweb}

\begin{grammarlens}[title={What you've built}]
One more word for telling, explaining and describing at B1.
\end{grammarlens}

\subsection*{Your turn}
\begin{itemize}
\raggedright
  \item \emph{Say it:} \emph{patrouiller}
  \item \emph{Say it:} \emph{patrouiller}, once more
  \item \emph{Say it:} a sentence of your own with \emph{patrouiller}
  \item \emph{Recall:} say the French for the gendarme, then the French for the gendarmerie, then say \emph{patrouiller} again
\end{itemize}

\begin{tcolorbox}[breakable,colback=teal!4,colframe=teal!35!black,title={Before you move on}]
What does patrouiller mean? (\textquotedblleft{}To patrol\textquotedblright{}.) Say it once more.
\end{tcolorbox}
\section[la patrouille]{la patrouille — the patrol}

\label{lesson:FR-C291-patrouille}



\begin{quote}
Before the new one: say the French for the gendarmerie, then the French for to patrol.
\end{quote}

\subsection*{You'll want to know: la patrouille}
\textbf{la patrouille} — \textquotedblleft{}the patrol\textquotedblright{}.

\begin{quote}
Une patrouille passe dans la rue toutes les heures.
\end{quote}

\begin{quote}
\emph{A patrol goes down the street every hour.}
\end{quote}

\emph{En patrouille} — on patrol.

\begin{cousinweb}
\emph{patrouille} comes from Old French \emph{patrouille} (EtymDB), built on \emph{patrouiller} \textquotedblleft{}to tramp about\textquotedblright{}.
\end{cousinweb}

\begin{grammarlens}[title={What you've built}]
One more word for telling, explaining and describing at B1.
\end{grammarlens}

\subsection*{Your turn}
\begin{itemize}
\raggedright
  \item \emph{Say it:} \emph{la patrouille}
  \item \emph{Say it:} \emph{la patrouille}, once more
  \item \emph{Say it:} a sentence of your own with \emph{la patrouille}
  \item \emph{Recall:} say the French for the gendarmerie, then the French for to patrol, then say \emph{la patrouille} again
\end{itemize}

\begin{tcolorbox}[breakable,colback=teal!4,colframe=teal!35!black,title={Before you move on}]
What does la patrouille mean? (\textquotedblleft{}The patrol\textquotedblright{}.) Say it once more.
\end{tcolorbox}
\section[la brigade]{la brigade — the squad, the brigade}

\label{lesson:FR-C291-brigade}



\begin{quote}
Before the new one: say the French for to patrol, then the French for the patrol.
\end{quote}

\subsection*{You'll want to know: la brigade}
\textbf{la brigade} — \textquotedblleft{}the squad, the brigade\textquotedblright{}.

\begin{quote}
Une brigade de police est arrivée en dix minutes.
\end{quote}

\begin{quote}
\emph{A police squad arrived in ten minutes.}
\end{quote}

\emph{La brigade des stupéfiants} — the drug squad.

\begin{cousinweb}
\emph{brigade} was borrowed from Italian \emph{brigata} \textquotedblleft{}a troop, a company\textquotedblright{} (EtymDB), built on \emph{briga} \textquotedblleft{}strife, trouble\textquotedblright{}.
\end{cousinweb}

\begin{grammarlens}[title={What you've built}]
One more word for telling, explaining and describing at B1.
\end{grammarlens}

\subsection*{Your turn}
\begin{itemize}
\raggedright
  \item \emph{Say it:} \emph{la brigade}
  \item \emph{Say it:} \emph{la brigade}, once more
  \item \emph{Say it:} a sentence of your own with \emph{la brigade}
  \item \emph{Recall:} say the French for to patrol, then the French for the patrol, then say \emph{la brigade} again
\end{itemize}

\begin{tcolorbox}[breakable,colback=teal!4,colframe=teal!35!black,title={Before you move on}]
What does la brigade mean? (\textquotedblleft{}The squad, the brigade\textquotedblright{}.) Say it once more.
\end{tcolorbox}
\section[(dialogue)]{Retour sur les chapitres 284–291}

\label{lesson:FR-R291-revision-284-291}



\begin{quote}
Say the French for to reveal and to report.
\end{quote}

\subsection*{Your turn}
Say these aloud:

\begin{itemize}
\raggedright
  \item to reveal, to report, to deny, to recount, to comment on
  \item the rumour, the scoop, disinformation, censorship, the controversy
  \item the poll, the final count, the rise, the drop, the statistic
  \item to burgle, the burglar, the burglary, the forced entry, the loot
  \item the pickpocket, the scam, the swindler, to swindle, the fraud
  \item the attack, to attack, the attacker, to threaten, the threat
  \item the crime, the criminal, the murder, the offender, crime
  \item the gendarme, the gendarmerie, to patrol, the patrol, the squad
\end{itemize}

\begin{tcolorbox}[breakable,colback=teal!4,colframe=teal!35!black,title={Before you move on}]
To reveal? (\textbf{révéler}.) To report? (\textbf{rapporter}.) To deny? (\textbf{démentir}.) To recount? (\textbf{relater}.) To comment on? (\textbf{commenter}.) The rumour? (\textbf{la rumeur}.) The scoop? (\textbf{le scoop}.) Disinformation? (\textbf{la désinformation}.) Censorship? (\textbf{la censure}.) The controversy? (\textbf{la polémique}.) The poll? (\textbf{le sondage}.) The final count? (\textbf{le bilan}.) The rise? (\textbf{la hausse}.) The drop? (\textbf{la baisse}.) The statistic? (\textbf{la statistique}.) To burgle? (\textbf{cambrioler}.) The burglar? (\textbf{le cambrioleur}.) The burglary? (\textbf{le cambriolage}.) The forced entry? (\textbf{l'effraction}.) The loot? (\textbf{le butin}.) The pickpocket? (\textbf{le pickpocket}.) The scam? (\textbf{l'arnaque}.) The swindler? (\textbf{l'escroc}.) To swindle? (\textbf{escroquer}.) The fraud? (\textbf{l'escroquerie}.) The attack? (\textbf{l'agression}.) To attack? (\textbf{agresser}.) The attacker? (\textbf{l'agresseur}.) To threaten? (\textbf{menacer}.) The threat? (\textbf{la menace}.) The crime? (\textbf{le crime}.) The criminal? (\textbf{le criminel}.) The murder? (\textbf{le meurtre}.) The offender? (\textbf{le délinquant}.) Crime? (\textbf{la délinquance}.) The gendarme? (\textbf{le gendarme}.) The gendarmerie? (\textbf{la gendarmerie}.) To patrol? (\textbf{patrouiller}.) The patrol? (\textbf{la patrouille}.) The squad? (\textbf{la brigade}.)
\end{tcolorbox}
